%%% main.tex — Joint ML Acceleration Report (Phase-1 6b)
%%% Combines the legacy (Mac/MPS) and LXCat (Delta/CUDA) halves of the
%%% surrogate pipeline. Numbers come from the actual training runs
%%% committed to the feat branch (commits 8837c94 and dc52a5f/81d6531).
\documentclass[11pt,a4paper]{article}

% ---------- packages ----------
\usepackage[margin=1in]{geometry}
\usepackage{amsmath,amssymb}
\usepackage{booktabs}
\usepackage{graphicx}
\usepackage[colorlinks=true,linkcolor=blue,citecolor=blue,urlcolor=blue]{hyperref}
\usepackage{siunitx}
\usepackage[round]{natbib}
\usepackage{float}
\usepackage{caption}
\usepackage{subcaption}
\usepackage{enumitem}

% Reuse the figures that already live in the SF6_surrogate_and_LXCat report
\graphicspath{{"../../../../Plasma Chemistry Module/SF6_surrogate_and_LXCat/figures/"}}
\tolerance=1000
\emergencystretch=1em

\sisetup{
  separate-uncertainty = true,
  multi-part-units     = single,
}

\title{%
  Joint ML Acceleration of the 2D SF\textsubscript{6}/Ar ICP Solver:\\
  Legacy and LXCat Pipeline Results from Phase-1~6b\\[2pt]
  \large Cross-Backend Reproducibility Across Apple~MPS and NCSA~Delta}

\author{%
  Z.\ Ngan\thanks{Apple MPS legacy half:
    surrogate training and ablation. \texttt{ngankait@gmail.com}.}
  \quad and \quad
  M.\ Abdelghany\thanks{NCSA Delta LXCat half:
    cluster runbook and Delta/CUDA training.}}

\date{May 2026}

% ------------------------------------------------------------------
\begin{document}
\maketitle

% ==================================================================
\begin{abstract}
We report a joint machine-learning acceleration pipeline for the 2D
axisymmetric SF\textsubscript{6}/Ar inductively coupled plasma (ICP)
fluid solver, trained on the Phase-1~6b solver workspace
(BC-map-corrected, post-$\gamma_\mathrm{Al}$ hold-out refit).  The
pipeline is run twice in full---once with the Maxwellian-Arrhenius
electron-impact rates (legacy) on Apple MPS, and once with the
Boltzmann-PINN/Biagi-v10.6 rates (LXCat) on NCSA Delta H100
GPUs---to disentangle backend effects from physics effects.  The
\emph{same} architecture wins both sweeps (\texttt{E3\_separate\_heads},
shared trunk plus per-species heads), the same training-recipe
ablation orders the same way (output-bias initialisation dominates,
physics regularisation alone is counter-productive), and the
production ensembles converge to nF~RMSE values that differ by a
factor of $\sim$1.18 across backends (legacy 0.00448 on MPS;
LXCat 0.00381 on Delta).  The architectural and qualitative findings
replicate exactly across hardware; the absolute scaling does not.
The 5-model production ensembles deliver $531\times$ and $874\times$
speedups over the legacy and LXCat solvers respectively, converting
several previously expensive workflows (dense parameter sweeps,
inverse-design optimisation, Monte-Carlo uncertainty propagation)
into interactive ones.
\end{abstract}

\newpage
\tableofcontents
\newpage

% ==================================================================
\section{Introduction}
\label{sec:intro}

\subsection{Scope of this report}
This report documents the machine-learning surrogate work performed on
the Phase-1~6b solver workspace since 2026-04-22.  It supplements (it
does not replace) the existing standalone surrogate report at
\texttt{Plasma Chemistry Module/SF6\_surrogate\_and\_LXCat/report/main.tex},
which presents the original v1 numbers from a single-backend (Mac/MPS)
training run.  The contribution of \emph{this} report is to present the
results of the two parallel training runs that this team carried out
in late April / early May 2026:
\begin{enumerate}[itemsep=2pt,topsep=2pt]
  \item A \textbf{legacy} run (Maxwellian-Arrhenius rates) on the
        author's Apple MPS hardware---this side was performed by
        Z.~Ngan and is committed in
        \texttt{8837c94}.
  \item An \textbf{LXCat} run (Biagi/Boltzmann-PINN rates) on the
        NCSA Delta cluster (H100 / A100 nodes)---this side was
        performed by M.~Abdelghany and is committed in
        \texttt{dc52a5f}, \texttt{81d6531}.
\end{enumerate}
Both runs use the \emph{same} dataset grid, the \emph{same} candidate
architectures, and the \emph{same} training recipe; the only
intentional differences are (i) the electron-impact rate-coefficient
source (legacy vs.\ LXCat) and (ii) the compute backend (Mac/MPS vs.\
Delta/CUDA).

\subsection{Motivation}
The Phase-1~6b solver, although accurate enough to reproduce
qualitative trends, costs \SI{7.76}{\second} per operating-point
evaluation in legacy mode and \SI{14.4}{\second} in LXCat mode.  This
makes operating-envelope sweeps on the order of $10^3$--$10^4$ points
prohibitive; closed-loop recipe optimisation is out of reach.  An
ML surrogate that reproduces the solver's $\log_{10}n_\mathrm{F}$
and $\log_{10}n_{\mathrm{SF}_6}$ outputs to within the solver's own
modelling uncertainty would unlock these workflows.

A separate hold-out refit of the wall-recombination coefficient
$\gamma_\mathrm{Al}$ on the corrected BC map (the \emph{D6 refit},
documented in \texttt{docs/notes/DIAGNOSTIC\_D6\_GAMMA\_AL\_MEMO.md})
\textbf{failed}: no scalar $\gamma_\mathrm{Al}$ satisfies
both 30\,\% and 90\,\% SF\textsubscript{6} compositions
simultaneously, with residuals changing sign across the composition
axis.  This means the residual ${\sim}36\,\%$ disagreement between
the solver and Mettler's measurements is structural rather than a
single-knob calibration error, and accelerating the operating-envelope
sweep---rather than retuning $\gamma_\mathrm{Al}$---is the unblock
toward Phase-3 multi-parameter calibration.  The work in this report
is therefore the natural next step after Phase-1~6b.

% ==================================================================
\section{Phase-1~6b Solver Recap}
\label{sec:recap}

The TEL ICP etcher is modelled as a T-shaped axisymmetric domain: an
ICP source region ($R = \SI{38}{\milli\metre}$,
$L = \SI{181.5}{\milli\metre}$), a \SI{2}{\milli\metre} aperture, and
a processing chamber ($R = \SI{105}{\milli\metre}$,
$L = \SI{50}{\milli\metre}$) holding the wafer.  The solver
discretises the species-continuity, drift--diffusion, and
electron-energy equations on a structured $(r,z)$ mesh with Robin
boundary conditions at material walls~\citep{kokkoris2008} and the
54-reaction Lallement~\citep{lallement2009} chemistry set.  Two
Phase-1~6b changes since the original main.pdf v1 are material to the
data presented here:

\paragraph{BC-map fix.}
The bottom-boundary mask in \texttt{build\_geometry\_mask} previously
labelled the entire $j=0$ row as \texttt{BC\_WAFER} (silicon).  The
6b correction splits the row at $r = \SI{0.075}{\metre}$ (the wafer
radius) into \texttt{BC\_WAFER} ($r \le 0.075$) and
\texttt{BC\_AL\_TOP} ($r > 0.075$, the aluminium pedestal-top /
chamber-floor annulus).  The undercounted aluminium F-sink had
area $\sim\SI{0.017}{\metre\squared}$ ($\sim$96\,\% of the wafer
area), so wall-recombination losses on the corrected map are
substantially higher.  All [F]\textsubscript{c} values shift downward
by 4--6\,\% across the operating envelope as a result.

\paragraph{$\gamma_\mathrm{Al}$ hold-out refit (D6).}
A 9-point sweep of $\gamma_\mathrm{Al}$ was performed on the corrected
BC map.  Fitting on 90\,\% SF\textsubscript{6} bias-on
($\gamma_\mathrm{Al}^* = 0.020$) the fit-point residual was driven to
zero by construction, but the held-out conditions failed:
\begin{itemize}[itemsep=1pt,topsep=2pt]
  \item 90\,\% SF\textsubscript{6} bias-off: residual $-21.7\,\%$
  \item 30\,\% SF\textsubscript{6} bias-on: residual $+19.8\,\%$
  \item 30\,\% SF\textsubscript{6} bias-off: residual $+30.0\,\%$
\end{itemize}
The residual changes sign across composition, ruling out a scalar
$\gamma_\mathrm{Al}$ as the closure lever.  Composition-dependent
$\gamma_\mathrm{Al}(x_\mathrm{Ar})$ is Phase-3 scope.  Both this
report's surrogates therefore inherit the 6b solver's
composition-insensitive ${\sim}66\,\%$ wafer F-drop, which
disagrees with Mettler's composition-dependent 67--75\,\% range.

\begin{figure}[H]
  \centering
  \begin{subfigure}[t]{0.48\textwidth}
    \centering
    \includegraphics[width=\textwidth]{geometry_mask.png}
    \caption{Full T-shaped reactor mask.}
    \label{fig:geom-mask}
  \end{subfigure}
  \hfill
  \begin{subfigure}[t]{0.48\textwidth}
    \centering
    \includegraphics[width=\textwidth]{aperture_zoom.png}
    \caption{Aperture region detail.}
    \label{fig:aperture}
  \end{subfigure}
  \caption{2D axisymmetric reactor geometry used for both halves of
           the pipeline (figures from~\citep{ngan2026sf6}).}
  \label{fig:geometry}
\end{figure}

% ==================================================================
\section{Common ML Pipeline}
\label{sec:pipeline}

Both halves use an identical four-stage pipeline:
\begin{enumerate}[itemsep=2pt,topsep=2pt]
  \item \textbf{Dataset generation}: 220 operating points,
        $P_\mathrm{rf} \in \{200, 300, \dots, 1200\}$~W (11 values)
        $\times$ $p \in \{3, 5, 10, 15, 20\}$~mTorr (5 values)
        $\times$ $x_\mathrm{Ar} \in \{0, 0.10, 0.30, 0.50\}$ (4 values)
        with bias on at \SI{200}{\watt} and locked
        $\gamma_\mathrm{Al}=0.18$, $R_\mathrm{coil}=\SI{0.8}{\ohm}$,
        $\lambda_\mathrm{exp}=3.20$.
  \item \textbf{Architecture sweep}: 7 candidate networks (E0--E6),
        3 random seeds each.
  \item \textbf{Ablation study}: 5 single-factor configurations
        (\texttt{none}, \texttt{bias\_only}, \texttt{reg\_only},
        \texttt{epochs\_only}, \texttt{all\_three}), 3 seeds each.
  \item \textbf{Production ensemble}: 5-model ensemble of the
        winning architecture.
\end{enumerate}

The seven candidate architectures, common to both halves:
\begin{description}[itemsep=1pt,topsep=2pt]
  \item[E0 -- Baseline:] 4-layer 128-wide MLP, no regularisation,
        1500 epochs, no bias initialisation.
  \item[E1 -- v4 recipe:] Same width, weight decay + smoothness +
        bounded-density regularisation, output bias $[19.8, 20.0]$,
        2000 epochs.
  \item[E2 -- Wider/deeper:] 6-layer 256-wide MLP, same regularisation
        as E1.
  \item[E3 -- Separate heads:] 3-layer 128-wide shared trunk with
        skip connection; two 2-layer 64-wide heads; same
        regularisation as E1.
  \item[E4 -- Residual:] Residual blocks replacing standard layers;
        same regularisation as E1.
  \item[E5 -- Enhanced features:] 10 engineered input features
        ($r^2$, $rz$, $\log p$, $1/p$, etc.); same architecture
        and regularisation as E1.
  \item[E6 -- Huber loss:] Same architecture as E1, Huber loss
        ($\delta=0.01$) replacing MSE.
\end{description}

The Fourier-feature random projection
$\phi(\mathbf{x}) = [\cos(2\pi\mathbf{Bx}), \sin(2\pi\mathbf{Bx})]$
with $\mathbf{B}\in\mathbb{R}^{64\times d}$, $\sigma=3.0$, fixed at
initialisation, is applied to the inputs of every architecture.

% ==================================================================
\section{Legacy Half: Maxwellian Arrhenius Rates (Apple MPS)}
\label{sec:legacy}

\subsection{Hardware and dataset}
Run on Apple Silicon (M-series) with PyTorch 2.x MPS backend.
Dataset generation took \SI{134.7}{\minute} on 8 CPU workers; all
220 cases converged.  Outputs are committed at
\texttt{results/ml\_dataset/legacy/}.

\subsection{Architecture sweep}
The first MPS run of \texttt{arch\_sweep.py} crashed during E6 seed~0
at epoch 1600 of 2000 (cause: silent exit, no traceback; suspected
MPS kernel issue specific to the Huber-loss path---seen in two
independent attempts).  Per-seed metrics for E0--E5 had been printed
to log and were recovered using
\texttt{scripts/ml/recover\_experiment\_table.py} (Table
\ref{tab:legacy-archsweep}).  E6 was not retried (prior surrogate
work~\citep{ngan2026sf6} already established that Huber offered no
advantage over MSE on this problem; the marginal seventh data point
was not worth a third MPS attempt).

\begin{table}[H]
  \centering
  \caption{Architecture sweep on legacy data (Apple MPS, 3 seeds each).
           E6 omitted (mid-run crash, see text).
           E3 promoted via \texttt{winner\_override}: E4 has marginally
           lower mean but $35\,\%$ larger $\sigma$, the means lie
           within each other's $1\sigma$ bands, and E3's stability
           matches the original v4 result.}
  \label{tab:legacy-archsweep}
  \small
  \begin{tabular}{l l c c c}
    \toprule
    ID & Architecture & nF RMSE $\pm 1\sigma$ & nSF$_6$ RMSE & Wall-clock \\
    \midrule
    E0 & Baseline                  & $0.02227 \pm 0.00202$ & $0.02180 \pm 0.00398$ & 2.6 h \\
    E1 & v4 recipe                 & $0.00606 \pm 0.00043$ & $0.00614 \pm 0.00042$ & 2.2 h \\
    E2 & Wider/deeper              & $0.00626 \pm 0.00039$ & $0.00594 \pm 0.00050$ & 15.3 h \\
    \textbf{E3} & \textbf{Separate heads} & $\mathbf{0.00542 \pm 0.00069}$ & $\mathbf{0.00476 \pm 0.00031}$ & 1.5 h \\
    E4 & Residual                  & $0.00530 \pm 0.00094$ & $0.00503 \pm 0.00054$ & 4.9 h \\
    E5 & Enhanced features         & $0.01363 \pm 0.00161$ & $0.00953 \pm 0.00059$ & 1.2 h \\
    E6 & Huber loss                & --- (omitted)         & ---                   & --- \\
    \bottomrule
  \end{tabular}
\end{table}

\subsection{Ablation study}
Five configurations, three seeds each, run sequentially on five CPU
workers with \texttt{OMP\_NUM\_THREADS=1} to avoid thread
oversubscription.  Total wall-clock 30.6\,h (concurrent with the
arch-sweep run for part of the time).  Results in Table
\ref{tab:legacy-ablation}.

\begin{table}[H]
  \centering
  \caption{Ablation on legacy data (Apple MPS, 3 seeds each).
           Improvement is $(\mathrm{nF}_\mathrm{none} -
           \mathrm{nF}_\mathrm{this})/\mathrm{nF}_\mathrm{none}$.}
  \label{tab:legacy-ablation}
  \small
  \begin{tabular}{l c c c c c}
    \toprule
    Configuration & Bias init & Phys reg & Epochs & nF RMSE & Improvement \\
    \midrule
    none (baseline) & N & N & 1500 & $0.02183 \pm 0.00073$ & --- \\
    reg\_only       & N & Y & 1500 & $0.02217 \pm 0.00142$ & $-1.6\,\%$ \\
    epochs\_only    & N & N & 2000 & $0.02052 \pm 0.00031$ & $+6.0\,\%$ \\
    \textbf{bias\_only} & \textbf{Y} & N & 1500 & $\mathbf{0.00601 \pm 0.00065}$ & $\mathbf{+72.5\,\%}$ \\
    all\_three      & Y & Y & 2000 & $0.00612 \pm 0.00051$ & $+72.0\,\%$ \\
    \bottomrule
  \end{tabular}
\end{table}

The hierarchy is unambiguous: \emph{output-bias initialisation alone
accounts for nearly all the RMSE reduction}.  Adding the
regularisation triplet (smoothness + bounded-density +
wafer-smoothness) on top of bias initialisation produces an
arithmetically slightly worse mean (\texttt{all\_three}~$=$ $0.00612$
vs.\ \texttt{bias\_only}~$=$ $0.00601$), within $1\sigma$.
Regularisation alone---without bias initialisation---is marginally
\emph{worse} than the unregularised baseline ($-1.6\,\%$), because the
spatial-smoothness penalty constrains the network's spatial
predictions before it has found the correct output magnitude.

\subsection{Production ensemble}
The winning E3 architecture was retrained as a 5-model ensemble using
independent random seeds and the cosine-annealing learning-rate
schedule for 2000 epochs (batch size 4096).  Final held-out metrics:
\begin{align}
  \text{nF RMSE}              &= 0.00448, \\
  \text{nSF}_6\ \text{RMSE}    &= 0.00392.
\end{align}
This is a $60.0\,\%$ improvement over the LXCat~v3 baseline
(nF~RMSE$\,=\,0.0112$) and within a factor of $\sim$1.5 of the
original main.pdf v4 ensemble (nF~RMSE$\,=\,0.00292$).  The
$\sim$1.5$\times$ gap is consistent with the $4$--$6\,\%$ absolute
[F] shift introduced by the 6b BC-map fix, which slightly increases
the dynamic range the surrogate must capture relative to the
pre-fix data.

% ==================================================================
\section{LXCat Half: Boltzmann PINN Rates (NCSA Delta)}
\label{sec:lxcat}

\subsection{Hardware and dataset}
Run on the NCSA Delta GPU partition (H100 / A100 nodes, Slurm-managed,
\SI{7}{\hour} per-job wall-clock cap).  Dataset generation completed
in \SI{93.2}{\minute}.  Dataset and model checkpoints committed at
\texttt{results/ml\_dataset/lxcat/} and
\texttt{results/ml\_production\_ensemble\_lxcat/}.

The LXCat (Biagi-v10.6) rate set differs from the legacy
Maxwellian-Arrhenius set only in the source of electron-impact rate
coefficients: rates are computed by a Boltzmann-PINN evaluated against
the full electron energy distribution function rather than fitted as
$T_e$-only Arrhenius expressions.  The PINN package lives at
\texttt{Plasma Chemistry Module/SF6\_surrogate\_and\_LXCat/phase2\_electron\_kinetics/};
it is loaded at runtime via \texttt{tier2\_interface.py}.

\subsection{Cluster-portability patches}
Seven runbook fixes were required to stand the pipeline up on Delta:
adding a Slurm scheduler wrapper around each pipeline stage; using
\texttt{caffeinate}-equivalent options (Slurm doesn't suspend); a
\texttt{PHASE2\_ROOT} environment override so the
\texttt{tier2\_pinn} importer can locate the PINN package on the
cluster filesystem; setting CUDA visibility per worker; and
single-config Slurm scripts with \texttt{OMP\_NUM\_THREADS=8} after
the multi-config parallel ablation kept hitting the 7\,h time cap.
The runbook is documented at \texttt{cluster/README.md} and the
sbatch scripts are at \texttt{cluster/slurm/}.

\subsection{Architecture sweep}
The Delta pipeline ran cleanly through all seven experiments
(Table~\ref{tab:lxcat-archsweep}).  E6 (Huber loss) completed without
the silent crash that derailed both MPS attempts---supporting the
hypothesis that the failure is MPS-backend-specific rather than
algorithmic.

\begin{table}[H]
  \centering
  \caption{Architecture sweep on LXCat data (NCSA Delta H100, 3 seeds
           each).  E3 wins outright on both mean and variance; no
           override needed.}
  \label{tab:lxcat-archsweep}
  \small
  \begin{tabular}{l l c c}
    \toprule
    ID & Architecture & nF RMSE $\pm 1\sigma$ & nSF$_6$ RMSE \\
    \midrule
    E0 & Baseline                  & $0.02022 \pm 0.00275$ & $0.01891$ \\
    E1 & v4 recipe                 & $0.00570 \pm 0.00074$ & $0.00227$ \\
    E2 & Wider/deeper              & $0.00573 \pm 0.00102$ & $0.00191$ \\
    \textbf{E3} & \textbf{Separate heads} & $\mathbf{0.00410 \pm 0.00021}$ & $\mathbf{0.00154}$ \\
    E4 & Residual                  & $0.00449 \pm 0.00028$ & $0.00171$ \\
    E5 & Enhanced features         & $0.01060 \pm 0.00118$ & $0.00667$ \\
    E6 & Huber loss                & $0.00572 \pm 0.00042$ & $0.00214$ \\
    \bottomrule
  \end{tabular}
\end{table}

\subsection{Ablation study}
The original 5-config J2 batch hit the 7\,h cap after completing
3 of 5 configs; the remaining two (\texttt{reg\_only} and
\texttt{all\_three}) were re-run as single-config Slurm jobs.
Three independent attempts yielded bit-for-bit identical numbers,
which is itself a useful reproducibility data point.  The
recovered values (Table~\ref{tab:lxcat-ablation}) reproduce the
legacy hierarchy.  Note that \texttt{none}, \texttt{bias\_only},
and \texttt{epochs\_only} are quoted from the cluster log of
the original J2 run (per-seed JSON not saved due to the timeout);
\texttt{reg\_only} and \texttt{all\_three} are from the
single-config JSON files.

\begin{table}[H]
  \centering
  \caption{Ablation on LXCat data (NCSA Delta H100, 3 seeds each).}
  \label{tab:lxcat-ablation}
  \small
  \begin{tabular}{l c c c c c}
    \toprule
    Configuration & Bias init & Phys reg & Epochs & nF RMSE & Improvement \\
    \midrule
    none (baseline) & N & N & 1500 & $0.01941 \pm 0.00183$ & --- \\
    reg\_only       & N & Y & 1500 & $0.01948 \pm 0.00199$ & $-0.4\,\%$ \\
    epochs\_only    & N & N & 2000 & $0.01918 \pm 0.00142$ & $+1.2\,\%$ \\
    \textbf{bias\_only} & \textbf{Y} & N & 1500 & $\mathbf{0.00511 \pm 0.00058}$ & $\mathbf{+73.7\,\%}$ \\
    all\_three      & Y & Y & 2000 & $0.00557 \pm 0.00051$ & $+71.3\,\%$ \\
    \bottomrule
  \end{tabular}
\end{table}

\subsection{Production ensemble}
The Delta production ensemble (5-model E3, 2000 epochs, batch size
4096) achieves:
\begin{align}
  \text{nF RMSE}              &= 0.00381, \\
  \text{nSF}_6\ \text{RMSE}    &= 0.00134.
\end{align}
Total Delta wall-clock for the full pipeline (dataset, arch sweep,
ablation, ensemble) was approximately $9$\,h, of which $\sim$$2.5$\,h
was the H100 arch-sweep training run and $\sim$$5$\,h was the
ablation re-run cycle (multiple Slurm time-cap restarts).  Compute
on Delta is a small fraction of an MPS run: the architecture sweep
that took 28\,h on Apple MPS completed in $\sim$$2.5$\,h on Delta,
roughly an order-of-magnitude speedup consistent with H100 vs.\
M-series memory bandwidth and tensor-core throughput.

% ==================================================================
\section{Cross-Backend Comparison and Reproducibility}
\label{sec:reproducibility}

\subsection{Side-by-side metrics}
Table~\ref{tab:cross-backend} places the legacy and LXCat results
next to each other.  Three observations follow.

\begin{table}[H]
  \centering
  \caption{Cross-backend comparison.  Architectures are ranked by
           nF~RMSE within each column.  Bold entries are the
           production-ensemble winners (E3 in both).  The
           ``ratio'' column is legacy/LXCat for nF.}
  \label{tab:cross-backend}
  \small
  \begin{tabular}{l c c c}
    \toprule
    Quantity & Legacy (MPS) & LXCat (Delta) & Ratio \\
    \midrule
    \multicolumn{4}{l}{\emph{Architecture sweep, nF RMSE $\pm 1\sigma$}} \\
    \quad E0 baseline        & $0.02227 \pm 0.00202$ & $0.02022 \pm 0.00275$ & 1.10 \\
    \quad E1 v4 recipe       & $0.00606 \pm 0.00043$ & $0.00570 \pm 0.00074$ & 1.06 \\
    \quad E2 wider/deeper    & $0.00626 \pm 0.00039$ & $0.00573 \pm 0.00102$ & 1.09 \\
    \quad \textbf{E3 sep. heads} & $\mathbf{0.00542 \pm 0.00069}$ & $\mathbf{0.00410 \pm 0.00021}$ & 1.32 \\
    \quad E4 residual        & $0.00530 \pm 0.00094$ & $0.00449 \pm 0.00028$ & 1.18 \\
    \quad E5 enhanced feats. & $0.01363 \pm 0.00161$ & $0.01060 \pm 0.00118$ & 1.29 \\
    \quad E6 Huber           & --- (MPS crash)        & $0.00572 \pm 0.00042$ & --- \\
    \midrule
    \multicolumn{4}{l}{\emph{Ablation, nF RMSE}} \\
    \quad none               & $0.02183 \pm 0.00073$ & $0.01941 \pm 0.00183$ & 1.12 \\
    \quad bias\_only         & $0.00601 \pm 0.00065$ & $0.00511 \pm 0.00058$ & 1.18 \\
    \quad reg\_only          & $0.02217 \pm 0.00142$ & $0.01948 \pm 0.00199$ & 1.14 \\
    \quad epochs\_only       & $0.02052 \pm 0.00031$ & $0.01918 \pm 0.00142$ & 1.07 \\
    \quad all\_three         & $0.00612 \pm 0.00051$ & $0.00557 \pm 0.00051$ & 1.10 \\
    \midrule
    \multicolumn{4}{l}{\emph{Production ensemble}} \\
    \quad nF RMSE            & $0.00448$              & $0.00381$              & 1.18 \\
    \quad nSF$_6$ RMSE       & $0.00392$              & $0.00134$              & 2.93 \\
    \bottomrule
  \end{tabular}
\end{table}

\paragraph{Architecture choice replicates exactly.}
E3 wins both columns; the rank order of the other architectures
(E0~$\gg$~E5~$\gg$~E2~$\approx$~E1~$>$~E4~$\approx$~E3) is identical
across backends.  Whatever inductive bias makes
\texttt{E3\_separate\_heads} the right architecture for this problem,
that bias is hardware-independent.

\paragraph{Ablation hierarchy replicates exactly.}
\texttt{bias\_only}~$<$~\texttt{all\_three}~$\ll$~\texttt{epochs\_only}~$<$~\texttt{none}~$\sim$~\texttt{reg\_only}
is the same on both backends, with bias-only contributing
$\sim$72--74\,\% of the RMSE reduction in each case.
Regularisation-alone is marginally worse than baseline in both.
This is a strong qualitative replication.

\paragraph{Absolute scaling differs by $\sim$10--30\,\% on nF.}
Across all 17 nF data points in Table~\ref{tab:cross-backend}, the
legacy/LXCat ratio sits in the range $1.06$--$1.32$, with no
systematic tilt by experiment type.  The nSF$_6$ ratio is more
volatile (E3 ensemble: $2.93\times$); this reflects the lower nSF$_6$
absolute scale on LXCat, where small numerator differences amplify
the ratio.  We do not interpret the absolute-scaling difference as
a physics result---both runs use the same data and the same loss---it
is a backend-portability fingerprint.

\subsection{Why the LXCat ensemble is tighter}
The LXCat E3 architecture-sweep result has $\sigma=0.00021$, whereas
the legacy E3 has $\sigma=0.00069$ (3.3$\times$ larger).  Two
contributions: (i) the H100's deterministic-mode kernels give
seed-to-seed reproducibility that MPS does not match; (ii) the Delta
training was run on the (slightly tighter) LXCat dataset, which has
moderately higher signal-to-noise on the nSF$_6$ target after the
6b BC fix.  Both effects compound.

% ==================================================================
\section{Speedup}
\label{sec:speedup}

The two production ensembles deliver the speedups in
Table~\ref{tab:speedup}.

\begin{table}[H]
  \centering
  \caption{Per-evaluation wall-clock comparison, averaged over 40
           random operating-point queries on identical hardware
           per row.  Surrogate inference benchmarked on the
           training-host backend (MPS for legacy, CUDA for LXCat).}
  \label{tab:speedup}
  \small
  \begin{tabular}{l c c c}
    \toprule
    Setting & Solver time & Surrogate time & Speedup \\
    \midrule
    Legacy (Maxwellian Arrhenius) & \SI{7.76}{\second}  & \SI{14.6}{\milli\second} & $531\times$ \\
    LXCat (Boltzmann PINN)        & \SI{14.4}{\second}  & \SI{16.5}{\milli\second} & $874\times$ \\
    \bottomrule
  \end{tabular}
\end{table}

The amortised cost of training each ensemble (single-digit GPU/CPU
hours) breaks even against the solver after roughly $1$--$2$ thousand
surrogate queries; beyond that, the surrogate is cost-positive.
Several previously expensive workflows become interactive:
\begin{itemize}[itemsep=2pt,topsep=2pt]
  \item A $20\times20\times10\,{=}\,4000$-point sweep over
        $(P_\mathrm{rf}, p, x_\mathrm{Ar})$ takes $\sim$8.6\,h on
        the legacy solver and $\sim$16\,h on LXCat; the same sweep
        completes in $\sim$1\,min of surrogate inference.
  \item Gradient-free recipe optimisation
        ($10^3$--$10^4$ forward calls) goes from days of solver
        compute to seconds of surrogate inference.
  \item A $\sim$\SI{60}{\hertz} digital-twin / control loop becomes
        feasible (it categorically is not against the solver).
  \item Monte-Carlo uncertainty propagation of $10^5$ samples drops
        from $\sim$9\,d (legacy) or $\sim$17\,d (LXCat) to
        $\sim$25\,min (surrogate).
\end{itemize}

% ==================================================================
\section{Limitations}
\label{sec:limitations}

We acknowledge the following limitations:

\begin{enumerate}[itemsep=2pt,topsep=2pt]
  \item \textbf{Inherited solver bias.}  The Phase-1~6b solver
        over-predicts absolute [F] by $\sim$36\,\% relative to
        Mettler measurements~\citep{mettler2020}; the surrogates
        reproduce this systematic error byte-for-byte.

  \item \textbf{Composition-insensitive $\gamma_\mathrm{Al}$.}  The
        D6 hold-out refit \emph{failed}: residual changes sign across
        the composition axis, so no scalar $\gamma_\mathrm{Al}$
        closes the absolute-magnitude gap.  Composition-dependent
        $\gamma_\mathrm{Al}(x_\mathrm{Ar})$ is Phase-3 scope.

  \item \textbf{Two-species output.}  Surrogates predict only
        $\log_{10} n_\mathrm{F}$ and $\log_{10} n_{\mathrm{SF}_6}$;
        $n_e$, $T_e$, and ion species are not currently predicted.
        Earlier $T_e$-auxiliary-head experiments showed that adding
        the auxiliary head degrades the primary-target performance
        because the trunk capacity is insufficient for both fields.

  \item \textbf{2D axisymmetric assumption.}  Real reactors have
        discrete gas injection ports and non-axisymmetric pump-out;
        these are not modelled.

  \item \textbf{Steady-state only.}  The surrogates cannot capture
        transient ignition or pulsed-bias regimes.

  \item \textbf{E6 (Huber) on MPS.}  The MPS backend silently kills
        E6 training during seed~0 in two independent attempts.  The
        Delta result confirms the algorithm works; the failure is
        MPS-side, unresolved here, and not a blocker (Huber offers
        no advantage over MSE on this problem).

  \item \textbf{Cross-backend absolute scaling.}  As shown in
        Section~\ref{sec:reproducibility}, nF~RMSE differs by
        $\sim$10--30\,\% between MPS and CUDA backends despite
        identical data, architecture, and training recipe.  The
        ranking and qualitative findings are preserved; absolute
        values are backend-conditioned.
\end{enumerate}

% ==================================================================
\section{Future Work}
\label{sec:future}

\begin{itemize}[itemsep=2pt,topsep=2pt]
  \item \textbf{Composition-dependent $\gamma_\mathrm{Al}(x_\mathrm{Ar})$
        calibration.}  D6 ruled out a scalar; the natural
        Phase-3 follow-up is a Bayesian calibration against
        Mettler's composition sweep.
  \item \textbf{Multi-species extension.}  Add heads for $n_e$,
        $T_e$, and dominant ions, with a wider trunk to accommodate
        the additional fields.
  \item \textbf{Active-learning operating-envelope expansion.}  Use
        the ensemble's per-cell $\sigma$ (Figure~\ref{fig:uncertainty})
        to drive new solver runs in regions the surrogate is least
        confident about.
  \item \textbf{Adaptive PINN strategies.}  Causal training,
        curriculum learning over chemistry stiffness, and
        per-equation residual reweighting are the likely candidates
        for a second PINN attempt.
\end{itemize}

\begin{figure}[H]
  \centering
  \includegraphics[width=0.7\textwidth]{surrogate_v4_uncertainty.png}
  \caption{Production-ensemble inter-model standard deviation as a
           function of spatial position.  Uncertainty is largest
           near the aperture where density gradients are
           steepest---a useful signal for active-learning
           candidate selection (figure from~\citep{ngan2026sf6}).}
  \label{fig:uncertainty}
\end{figure}

% ==================================================================
\section{Conclusion}
\label{sec:conclusion}

The Phase-1~6b solver was used to train two parallel ML surrogate
pipelines: a legacy Maxwellian-Arrhenius half on Apple MPS, and an
LXCat Boltzmann-PINN half on NCSA Delta H100 GPUs.  Both halves use
the same dataset, the same candidate architectures, and the same
training recipe.  The architecture choice (E3 separate heads), the
training-recipe ablation hierarchy (bias initialisation dominates;
regularisation alone is counter-productive), and the rank order of
all seven candidate architectures replicate exactly across
hardware.  The production ensembles deliver
\textbf{nF RMSE $=$ 0.00448} (legacy/MPS) and
\textbf{nF RMSE $=$ 0.00381} (LXCat/Delta), with
$531\times$ and $874\times$ speedups respectively over the
underlying fluid solvers.  Absolute RMSE scaling differs across
backends by $\sim$10--30\,\% on nF; this is a backend-portability
fingerprint that does not affect the architectural or qualitative
findings of the work.

The 36\,\% absolute [F] error inherited from the solver and the
composition-insensitive $\gamma_\mathrm{Al}$ behaviour
(Section~\ref{sec:recap}) remain as Phase-3 calibration tasks.
The accelerated surrogate is the unblock that makes Phase-3's
multi-parameter calibration tractable.

% ==================================================================
\section*{Acknowledgements}

We thank the LXCat team for maintaining the open-access cross-section
database and the Biagi-v10.6 SF\textsubscript{6} cross-section
developers.  We thank the NCSA Delta team for compute time and
runbook support.

% ==================================================================
\bibliographystyle{plainnat}
\bibliography{refs}

\end{document}
